

\documentclass[11pt]{article}

% =================================================
% Packages
% =================================================

% Page layout
\usepackage[margin=1in]{geometry}

% Mathematical packages
\usepackage{amsmath,amssymb,amsthm}
\usepackage{mathtools}
\usepackage{bm}
\usepackage{comment}
% Algorithms
\usepackage[ruled,vlined]{algorithm2e}

% Lists and tables
\usepackage{enumitem}
\usepackage{booktabs}
\usepackage{array}

% Figures and geometric illustrations
\usepackage{graphicx}
\usepackage{float}
\usepackage{tikz}
\usetikzlibrary{arrows.meta,calc,positioning}

% Colors and code listings
\usepackage{xcolor}
\usepackage{verbatim}
\usepackage{listings}

% Hyperlinks
\usepackage[
    colorlinks=true,
    linkcolor=blue,
    urlcolor=blue,
    citecolor=blue
]{hyperref}

% =================================================
% Code Listing Style
% =================================================

\lstset{
    frame=single,
    breaklines=true,
    basicstyle=\small\ttfamily,
    keywordstyle=\color{blue},
    commentstyle=\color{green!60!black},
    stringstyle=\color{purple},
    numbers=left,
    numberstyle=\tiny\color{gray}
}

% =================================================
% Formatting
% =================================================

\setlength{\parindent}{0pt}
\setlength{\parskip}{0.3cm}

% =================================================
% Theorem Environments
% =================================================

\newtheorem{theorem}{Theorem}[section]
\newtheorem{definition}[theorem]{Definition}
\newtheorem{proposition}[theorem]{Proposition}
\newtheorem{lemma}[theorem]{Lemma}
\newtheorem{corollary}[theorem]{Corollary}
\newtheorem{remark}[theorem]{Remark}
\newtheorem{example}[theorem]{Example}

% =================================================
% Title
% =================================================

\title{
\textbf{
Principal Component Analysis and Its
Mathematical Justification Through the
Eckart--Young Theorem
}\\[0.3cm]
\large
A Step-by-Step Development from Geometric
Intuition to Optimal Low-Rank Approximation
}

\author{Biajid Chowdhury}

\date{\today}

% =================================================
% DOCUMENT BEGINS
% =================================================

\begin{document}

\maketitle

% =================================================
% ABSTRACT
% =================================================

% Insert the abstract here.




% =================================================
% MAIN CONTENT
% =================================================

% We will develop each section individually.





% =================================================
% ABSTRACT
% =================================================

\begin{abstract}

In recent years, Principal Component Analysis (PCA) has gained
enormous popularity in data science, machine learning, and
statistical analysis. Despite its widespread application,
however, its underlying mathematical justification is often
given surprisingly little attention, particularly in
introductory data science curricula. Instead, PCA is
frequently introduced as a black-box computational tool:
subtract the mean, compute the Singular Value Decomposition
(SVD), retain a few dominant components, and proceed with
the transformed data. Although this approach is convenient,
it leaves a fundamental mathematical question unanswered:
\textit{Why should these particular components provide the
best possible lower-dimensional representation of the
original data?}

Fortunately, with relatively modest mathematical
background and some geometric intuition, one can
appreciate the remarkable beauty of PCA and understand
why it works. Such an understanding not only removes
the mystery behind the algorithm but also allows
one to apply it more intelligently and effectively.

In this note, we develop Principal Component Analysis
step-by-step, beginning with the geometric interpretation
of data projection, variance, and Singular Value
Decomposition. We then explain how principal components
are obtained, how they transform the original data,
and how a lower-dimensional representation can be
used to reconstruct an approximation of the original
dataset. Finally, we establish the mathematical
optimality of this reconstruction through the
Eckart--Young--Mirsky theorem, examining both the
Frobenius and spectral norms. Our objective is to
bridge the gap between the computational recipe
commonly taught in data science and the elegant
linear algebra that justifies it.

\end{abstract}

\tableofcontents
\newpage



% =================================================
\section{Introduction}
% =================================================

In modern data science and machine learning, one frequently
encounters datasets involving hundreds, thousands, or even
millions of variables. Although such high-dimensional data
may contain valuable information, working directly with
all these variables is not always computationally efficient
or mathematically relevant. Furthermore, many of
these variables may be strongly correlated, meaning that
a considerable amount of the variation in the data can
often be represented using substantially fewer dimensions.

This naturally raises an interesting question:
\textit{Can we represent a high-dimensional dataset
using only a few carefully chosen directions, while
preserving as much of its original variation as possible?}

Principal Component Analysis (PCA) provides an elegant
mathematical answer to this question. The central idea
is remarkably simple: rather than describing observations
using the original coordinate axes, we construct a new
set of mutually orthogonal directions along which the
data exhibits the greatest variation. By retaining only
the most significant of these directions, we obtain
a lower-dimensional representation of the original
dataset.

At first glance, the computational procedure appears
almost deceptively simple. We begin by subtracting
the mean of each variable, perform a Singular Value
Decomposition (SVD) of the centered data matrix, retain
a few singular vectors corresponding to the largest
singular values, and project the observations onto
the subspace spanned by these vectors.

But why should this procedure work?

More precisely, why should the singular vectors
associated with the largest singular values identify
the directions of maximum variation? What exactly
happens to the original observations when we project
them onto these directions? And perhaps most
importantly, if we reconstruct the original dataset
from its lower-dimensional representation, can
we guarantee that the resulting approximation is
the best possible one?

These questions reveal that PCA is much more than
a convenient computational algorithm. Beneath its
apparently simple implementation lies a beautiful
connection among orthogonal projections, eigenvalue
problems, matrix norms, and low-rank matrix
approximation.

The mathematical foundation of this connection is
provided by two closely related optimization
principles. First, the principal directions maximize
the variance captured by an orthogonal projection
onto a subspace of prescribed dimension. Second,
the same subspace minimizes the squared reconstruction
error when the projected observations are mapped
back into the original feature space.

The second property is closely connected to the
celebrated Eckart--Young theorem, which establishes
the optimality of truncated Singular Value
Decomposition among all matrices of a prescribed
maximum rank. In particular, the theorem shows
that retaining the largest singular values produces
an approximation with the smallest possible error,
measured in either the Frobenius norm or the
spectral norm.

In this note, our objective is not merely to
present another recipe for performing PCA. Rather,
we wish to develop the mathematical ideas that
make the procedure work, beginning with familiar
concepts from elementary linear algebra and
gradually arriving at its rigorous justification.

We begin with a geometric discussion of
high-dimensional data, orthogonal projections,
and the meaning of variance along a given direction.
We then formulate the problem of finding directions
of maximum variance and establish its connection
to the eigenvectors of the covariance matrix.

Next, we revisit Singular Value Decomposition
from the perspective of PCA. We explain the
geometric significance of the left and right
singular vectors, the role of singular values,
and how the principal components arise naturally
from this decomposition.

Having established these ideas, we examine the
process of dimensionality reduction and
reconstruction. In particular, we investigate
what information is retained, what information
is discarded, and how the reconstruction error
can be measured mathematically.

Finally, we state and prove the Eckart--Young
theorem for both the Frobenius norm and the
spectral norm. These proofs allow us to
establish precisely why the truncated SVD
provides an optimal low-rank approximation,
rather than simply accepting its optimality
as a computational fact.

Throughout the discussion, we emphasize
geometric intuition alongside mathematical
rigor. Whenever possible, abstract arguments
will be accompanied by concrete examples,
and the theoretical results will ultimately
be verified through numerical experiments
using Python.

Our hope is that, by the end of this note,
PCA will no longer appear as an obscure
black-box tool, but rather as a natural
consequence of some of the most fundamental
and elegant ideas in scientific computation and numerical linear algebra.



% =================================================
\section{Prerequisite Mathematical Background}
% =================================================

The primary objective of this note is to develop
the mathematical justification of Principal Component
Analysis rather than to introduce the elementary
concepts of linear algebra from the beginning.
Accordingly, we assume that the reader has some
familiarity with the fundamental ideas ordinarily
covered in an undergraduate linear algebra course.

In particular, the reader should be comfortable
with the following concepts:

\begin{itemize}[leftmargin=1.5cm]

    \item \textbf{Vector spaces and subspaces:}
    Linear independence, basis, dimension, rank,
    null space, and the rank-nullity theorem.

    \item \textbf{Inner products and orthogonality:}
    The Euclidean inner product, orthonormal bases,
    orthogonal matrices, and orthogonal projections
    onto subspaces.

    \item \textbf{Eigenvalues and eigenvectors:}
    Eigenvalue decomposition, diagonalization,
    symmetric matrices, and the spectral theorem
    for real symmetric matrices.

    \item \textbf{Singular Value Decomposition (SVD):}
    The factorization
    \[
        A = U\Sigma V^T,
    \]
    including the meaning of singular values,
    left and right singular vectors, and their
    relationship to the eigenvalues and
    eigenvectors of $AA^T$ and $A^TA$.

    \item \textbf{Basic matrix operations:}
    Matrix multiplication, transposition,
    matrix-vector multiplication, and elementary
    properties of matrix norms.

\end{itemize}

Readers who are unfamiliar with these topics,
or who have studied them previously but would
benefit from a review, are encouraged to consult
Gilbert Strang's \textit{Linear
Algebra and Its Applications} before proceeding. In particular,
the discussions of orthogonality, projections,
eigenvalues, symmetric matrices, and Singular
Value Decomposition provide an excellent
foundation for understanding the material
developed here.

Nevertheless, we do not expect the reader
to be familiar with Principal Component
Analysis, the Eckart--Young theorem, or
the optimization problems associated with
low-rank approximation. These topics will
be developed carefully in the following
sections.

We will also introduce the Frobenius and
spectral norms explicitly and explain their
geometric significance before using them
in the proofs. Similarly, the necessary
statistical concepts, particularly variance
and covariance, will be reviewed in the
context of PCA.

Although some of our arguments will involve
optimization, no advanced background in
optimization theory is assumed. The required
results will be motivated and established
using elementary linear algebra whenever
possible.

Finally, we emphasize that this note is
intended to be read with pencil and paper.
The reader is encouraged to verify the
matrix identities, work through the examples,
and examine the geometric meaning of each
transformation. Our goal is not simply
to arrive at the correct formulas, but
to understand why those formulas must
take the form they do.









% =================================================
\section{A Brief Review of Singular Value
Decomposition}
% =================================================

Before proceeding to the geometric and mathematical
development of Principal Component Analysis,
it is worthwhile to briefly revisit the Singular
Value Decomposition (SVD) of a matrix.

We do not intend to provide a complete treatment
of SVD here, as this is ordinarily covered in
undergraduate linear algebra courses. Instead,
we shall recall its fundamental properties,
particularly the relationship between singular
vectors, the row space, the column space, and
orthogonal projections.

These ideas will play a central role in our
subsequent development of PCA.

% -------------------------------------------------
\subsection{Row Space and Column Space}
% -------------------------------------------------

Let $A\in\mathbb{R}^{m\times n}$ be an arbitrary
real matrix.

There are two important vector spaces naturally
associated with $A$.

First, consider the columns of $A$ as vectors
in $\mathbb{R}^m$. Their linear combinations
form the \textit{column space} of $A$, denoted by

\[
\operatorname{Col}(A)
=
\operatorname{span}\{a_1,a_2,\ldots,a_n\},
\]

where $a_j$ denotes the $j$th column of $A$.

In other words, the column space consists of
all vectors that can be obtained by taking
linear combinations of the columns of $A$.

Similarly, we may consider each row of $A$
as a vector in $\mathbb{R}^n$.
The linear combinations of these row vectors
form the \textit{row space} of $A$, denoted by

\[
\operatorname{Row}(A)
=
\operatorname{Col}(A^T).
\]

Thus, the column space is a subspace of
$\mathbb{R}^m$, whereas the row space is
a subspace of $\mathbb{R}^n$.

Although these two spaces generally belong
to different Euclidean spaces, an important
theorem from elementary linear algebra tells
us that they have the same dimension:

\[
\boxed{
\dim\operatorname{Col}(A)
=
\dim\operatorname{Row}(A)
=
\operatorname{rank}(A).
}
\]

If $\operatorname{rank}(A)=r$, both spaces
therefore admit bases consisting of
$r$ linearly independent vectors.

There is, however, no reason to insist
that these bases consist of the original
rows or columns of $A$.

Indeed, it is often much more convenient
to construct orthonormal bases for these
spaces. Singular Value Decomposition provides
precisely such bases, together with an
additional and rather remarkable structure.

% -------------------------------------------------
\subsection{The Singular Value Decomposition Theorem}
% -------------------------------------------------

\begin{theorem}[Singular Value Decomposition]

For every real matrix
$A\in\mathbb{R}^{m\times n}$,
there exist orthogonal matrices

\[
U\in\mathbb{R}^{m\times m},
\qquad
V\in\mathbb{R}^{n\times n},
\]

and a rectangular diagonal matrix

\[
\Sigma\in\mathbb{R}^{m\times n},
\]

such that

\[
\boxed{A=U\Sigma V^T.}
\]

The matrices $U$ and $V$ satisfy

\[
U^TU=UU^T=I_m,
\qquad
V^TV=VV^T=I_n.
\]

The diagonal entries of $\Sigma$ are
nonnegative real numbers, called the
\textit{singular values} of $A$.

If $\operatorname{rank}(A)=r$, the positive
singular values can be arranged in
nonincreasing order:

\[
\sigma_1\geq\sigma_2\geq
\cdots\geq\sigma_r>0.
\]

All remaining diagonal entries of
$\Sigma$ are zero.

\end{theorem}

The columns of $U$ are called the
\textit{left singular vectors}, while
the columns of $V$ are called the
\textit{right singular vectors}.

We write

\[
U=
\begin{bmatrix}
u_1 & u_2 & \cdots & u_m
\end{bmatrix},
\]

and

\[
V=
\begin{bmatrix}
v_1 & v_2 & \cdots & v_n
\end{bmatrix}.
\]

Since these two matrices are orthogonal,
their respective columns form
orthonormal bases of $\mathbb{R}^m$
and $\mathbb{R}^n$.

More importantly, the first $r$
singular vectors provide orthonormal
bases for the fundamental spaces
associated with $A$:

\[
\boxed{
\operatorname{Col}(A)
=
\operatorname{span}
\{u_1,\ldots,u_r\}
}
\]

and

\[
\boxed{
\operatorname{Row}(A)
=
\operatorname{span}
\{v_1,\ldots,v_r\}.
}
\]

Thus, SVD identifies orthonormal bases
for both the column space and the
row space of an arbitrary matrix.

% -------------------------------------------------
\subsection{How Are the Singular Vectors Obtained?}
% -------------------------------------------------

The singular values and singular vectors
are closely related to the eigenvalue
decompositions of $AA^T$ and $A^TA$.

Starting from

\[
A=U\Sigma V^T,
\]

we obtain

\[
AA^T
=
U\Sigma\Sigma^TU^T,
\]

and similarly,

\[
A^TA
=
V\Sigma^T\Sigma V^T.
\]

Therefore, the columns of $U$ are
eigenvectors of $AA^T$, while the
columns of $V$ are eigenvectors
of $A^TA$.

The nonzero eigenvalues of these
two symmetric matrices are identical,
and are given by

\[
\boxed{\lambda_i=\sigma_i^2.}
\]

Consequently, the positive singular
values are obtained by taking the
square roots of the positive eigenvalues
of either $AA^T$ or $A^TA$:

\[
\sigma_i=\sqrt{\lambda_i}.
\]

For each positive singular value,
the corresponding left and right
singular vectors satisfy

\[
\boxed{
Av_i=\sigma_i u_i,
\qquad
A^Tu_i=\sigma_i v_i.
}
\]

In particular, if $v_i$ is a normalized
eigenvector of $A^TA$ corresponding
to a positive eigenvalue $\sigma_i^2$,
then

\[
\boxed{
u_i=\frac{Av_i}{\sigma_i}.
}
\]

Conversely, if $u_i$ is known, we can
obtain the corresponding right singular
vector through

\[
\boxed{
v_i=\frac{A^Tu_i}{\sigma_i}.
}
\]

These formulas apply whenever
$\sigma_i>0$. The remaining singular
vectors are chosen to complete the
respective orthonormal bases.

We may therefore interpret the
singular value decomposition as a
carefully coordinated pair of
eigenvalue decompositions.

It is worth emphasizing that singular
vectors are not necessarily unique.
Repeated singular values allow
different choices of orthonormal bases
within the corresponding singular
subspaces. Nevertheless, the
decomposition always exists.

% -------------------------------------------------
\subsection{The Reduced Singular Value Decomposition}
% -------------------------------------------------

Suppose $\operatorname{rank}(A)=r$.

Instead of retaining all the columns
of $U$ and $V$, we may retain only
those corresponding to positive
singular values.

Define

\[
U_r=
\begin{bmatrix}
u_1 & \cdots & u_r
\end{bmatrix}
\in\mathbb{R}^{m\times r},
\]

\[
V_r=
\begin{bmatrix}
v_1 & \cdots & v_r
\end{bmatrix}
\in\mathbb{R}^{n\times r},
\]

and

\[
\Sigma_r
=
\operatorname{diag}
(\sigma_1,\ldots,\sigma_r).
\]

Then

\[
\boxed{
A=U_r\Sigma_rV_r^T.
}
\]

This is called the \textit{reduced}
or \textit{compact} singular value
decomposition.

Equivalently, we may express $A$
as a sum of rank-one matrices:

\[
\boxed{
A=\sum_{i=1}^{r}
\sigma_i u_i v_i^T.
}
\]

Each term $\sigma_i u_i v_i^T$
has rank one.

This representation will become
particularly useful when we discuss
low-rank approximation and the
Eckart--Young theorem.

% -------------------------------------------------
\subsection{The Geometric Meaning of SVD}
% -------------------------------------------------

Let us now consider the matrix $A$
as a linear transformation

\[
A:\mathbb{R}^n\longrightarrow
\mathbb{R}^m.
\]

For an arbitrary vector
$x\in\mathbb{R}^n$, we have

\[
Ax=U\Sigma V^Tx.
\]

The transformation can be understood
as three consecutive operations.

First, $V^T$ expresses the input
vector $x$ in the orthonormal
coordinate system determined by
the right singular vectors.

Next, $\Sigma$ scales the resulting
coordinates by the corresponding
singular values. Coordinates associated
with zero singular values are
eliminated.

Finally, $U$ expresses the resulting
vector in the output space
$\mathbb{R}^m$ using the left
singular vectors.

Thus, schematically,

\[
\boxed{
x
\xrightarrow{\;V^T\;}
V^Tx
\xrightarrow{\;\Sigma\;}
\Sigma V^Tx
\xrightarrow{\;U\;}
Ax.
}
\]

In particular, since

\[
Av_i=\sigma_i u_i,
\]

the transformation $A$ takes a
right singular direction $v_i$
and maps it to the corresponding
left singular direction $u_i$,
scaled by $\sigma_i$.

Hence, the positive singular values
describe the stretching factors
along the corresponding singular
directions.

Geometrically, SVD identifies
orthogonal input directions,
their associated stretching
factors, and the corresponding
orthogonal output directions.

% -------------------------------------------------
\subsection{SVD and Orthogonal Projections}
% -------------------------------------------------

We now return to the row space
and column space of $A$.

Recall that the columns of $V_r$
form an orthonormal basis of
$\operatorname{Row}(A)$.

Therefore, for any
$x\in\mathbb{R}^n$, its
orthogonal projection onto the
row space is

\[
\boxed{
P_{\operatorname{Row}(A)}x
=
V_rV_r^Tx.
}
\]

Similarly, since the columns
of $U_r$ form an orthonormal
basis of $\operatorname{Col}(A)$,
the orthogonal projection of
any $y\in\mathbb{R}^m$ onto
the column space is

\[
\boxed{
P_{\operatorname{Col}(A)}y
=
U_rU_r^Ty.
}
\]

An important distinction deserves
attention here.

The matrix $V_r^T$ does not itself
perform the orthogonal projection
onto the row space in the original
ambient coordinates. Rather,

\[
V_r^Tx
\]

provides the coordinates of the
row-space projection relative to
the orthonormal basis
$\{v_1,\ldots,v_r\}$.

Multiplying these coordinates by
$V_r$ reconstructs the projected
vector in $\mathbb{R}^n$:

\[
V_r(V_r^Tx)=V_rV_r^Tx.
\]

The same interpretation applies
to the column space:

\[
U_r^Ty
\]

provides the coordinates of the
column-space projection, whereas

\[
U_r(U_r^Ty)=U_rU_r^Ty
\]

produces the actual projected
vector in $\mathbb{R}^m$.

Thus, we obtain the important
identities

\[
\boxed{
\begin{aligned}
P_{\operatorname{Row}(A)}
&=V_rV_r^T,\\[2mm]
P_{\operatorname{Col}(A)}
&=U_rU_r^T.
\end{aligned}
}
\]

These formulas are fundamental
consequences of the geometry
of SVD.

% -------------------------------------------------
\subsection{From the Row Space to the Column Space}
% -------------------------------------------------

There is another useful way to
understand the transformation $A$.

The Fundamental Theorem of Linear
Algebra tells us that

\[
\mathbb{R}^n
=
\operatorname{Row}(A)
\oplus
\operatorname{Null}(A),
\]

where the two subspaces are orthogonal.

Thus, any input vector
$x\in\mathbb{R}^n$ can be
decomposed uniquely as

\[
x=x_{\mathrm{row}}+
x_{\mathrm{null}},
\]

where

\[
x_{\mathrm{row}}
\in\operatorname{Row}(A),
\qquad
x_{\mathrm{null}}
\in\operatorname{Null}(A).
\]

Applying $A$, we obtain

\[
Ax
=
Ax_{\mathrm{row}}
+
Ax_{\mathrm{null}}.
\]

But $Ax_{\mathrm{null}}=0$.
Consequently,

\[
\boxed{
Ax=Ax_{\mathrm{row}}.
}
\]

Therefore, the action of $A$
depends only on the component
of $x$ lying in its row space.

The null-space component contributes
nothing to the output.

Furthermore, the resulting vector
necessarily belongs to the column
space of $A$.

In this sense, $A$ takes the relevant
information from its row space,
stretches it along the singular
directions, and maps it into
its column space.

Indeed, for the reduced SVD,

\[
Ax=U_r\Sigma_rV_r^Tx.
\]

The matrix $V_r^T$ extracts
the row-space coordinates,
$\Sigma_r$ scales them,
and $U_r$ constructs the
corresponding vector in
the column space.

This provides a particularly
useful geometric interpretation
of the singular value decomposition.

% -------------------------------------------------
\subsection{Why Is This Relevant to PCA?}
% -------------------------------------------------

When $A$ represents a dataset
with features arranged along
the rows and observations
along the columns, each
observation is a vector in
$\mathbb{R}^m$.

Thus, the left singular vectors
$u_i$ belong to the same
feature space as the
original observations.

As we shall establish in
the subsequent sections,
these vectors identify the
principal directions along
which the centered data
exhibits variation.

If we retain only the first
$k$ left singular vectors,
we obtain

\[
U_k=
\begin{bmatrix}
u_1 & \cdots & u_k
\end{bmatrix}.
\]

The matrix

\[
\boxed{Z=U_k^TA}
\]

then expresses each observation
in terms of these $k$
orthonormal directions.

We may reconstruct an
approximation in the
original feature space by

\[
\boxed{
\widehat{A}
=
U_kZ
=
U_kU_k^TA.
}
\]

At this stage, we have only
described the algebraic
meaning of this operation.

We have not yet established
why the first $k$ singular
vectors should be preferred
over other possible directions.

That question forms the
mathematical heart of PCA,
and we shall address it
through geometric intuition,
variance maximization, and
eventually the Eckart--Young
theorem.



\begin{comment}
    
% =================================================
\section{Geometric Motivation for Principal
Component Analysis}
% =================================================
The fundamental idea behind PCA is that a collection of
high-dimensional observations may exhibit most
of its variation along only a few directions.
If we can identify these directions, we may be
able to describe the dataset using substantially
fewer coordinates without losing too much of
its original structure.

But what exactly do we mean by a direction
that captures the variation of a dataset?
And how can we determine which direction
is the most important?

We begin with a simple geometric interpretation.

% -------------------------------------------------
\subsection{Representing a Dataset Geometrically}
% -------------------------------------------------

Suppose we have $n$ observations, each described
by $m$ features. We represent the dataset by

\[
X =
\begin{bmatrix}
| & | & & |\\
x_1 & x_2 & \cdots & x_n\\
| & | & & |
\end{bmatrix}
\in \mathbb{R}^{m\times n},
\]

where each column $x_j\in\mathbb{R}^m$
represents one observation, and each row
corresponds to a particular feature.

For example, if our observations are individuals,
the rows might represent age, income, height,
or years of education. Each individual is
therefore represented by a point in an
$m$-dimensional Euclidean space.

When $m=2$, the observations form a collection
of points in the plane. When $m=3$, they form
a cloud of points in three-dimensional space.

For larger $m$, direct visualization becomes
difficult, but the underlying geometry
remains unchanged.

Now imagine that this cloud of points is
not distributed uniformly in every direction.
Instead, it may be elongated along some
particular direction.

Intuitively, that direction should contain
more information about how the observations
differ from one another than a direction
along which the points are tightly clustered.

This observation provides the initial
motivation for PCA.

The idea is to identify a new system of
orthogonal coordinate axes that reflects
the dominant directions of variation
in the dataset.

Rather than describing each observation
using all the original features, we may
then describe it using only a few
coordinates along these new axes.

But how does PCA accomplish this?
Let us first examine the computational
procedure before investigating why
it is mathematically justified.

% -------------------------------------------------
\subsection{The Principal Component Analysis
Procedure}
% -------------------------------------------------

The basic PCA procedure can be summarized
in a few relatively simple steps.

Although different textbooks may arrange
the data matrix differently, we will
consistently assume that features occupy
the rows and observations occupy the
columns.

\subsubsection*{Step 1: Center the Data}

We begin by computing the mean
observation vector

\[
\bar{x}=
\frac{1}{n}\sum_{j=1}^{n}x_j.
\]

Each observation is then centered
by subtracting this mean:

\[
a_j=x_j-\bar{x}.
\]

The resulting centered data matrix is

\[
\boxed{
A=X-\bar{x}\mathbf{1}_n^T,
}
\]

where $\mathbf{1}_n$ is the vector
of $n$ ones.

Notice that

\[
\sum_{j=1}^{n}a_j=0.
\]

Consequently, the center of the
data cloud has been translated
to the origin.

This step is important because
our objective is to measure
the variation of observations
around their mean, rather than
their distance from an arbitrarily
chosen origin.

From now on, unless stated
otherwise, $A$ denotes the
centered data matrix.

\subsubsection*{Step 2: Compute the SVD}

We compute the Singular Value
Decomposition of $A$:

\[
\boxed{
A=U\Sigma V^T.
}
\]

Recall from Section 3 that
the columns of $U$ are
the left singular vectors
and the columns of $V$
are the right singular vectors.

The positive singular values
are arranged in decreasing order:

\[
\sigma_1\geq\sigma_2
\geq\cdots\geq\sigma_r>0,
\]

where $r=\operatorname{rank}(A)$.

Since $A$ is an $m\times n$
matrix, its left singular
vectors belong to $\mathbb{R}^m$,
the original feature space.

Thus, these vectors provide
candidate directions along
which we may describe the
observations.

But which of these directions
should we retain?

The PCA procedure selects
those corresponding to the
largest singular values.

We shall shortly investigate
why this choice is justified.

\subsubsection*{Step 3: Select the Principal
Directions}

Suppose we wish to represent
the original $m$-dimensional
observations using only
$k$ coordinates, where

\[
1\leq k\leq r.
\]

We retain the first $k$
columns of $U$:

\[
\boxed{
U_k=
\begin{bmatrix}
u_1 & u_2 & \cdots & u_k
\end{bmatrix}
\in\mathbb{R}^{m\times k}.
}
\]

These vectors form an
orthonormal basis for a
$k$-dimensional subspace
of the original feature space.

We call this the
\textit{principal subspace}.

At this stage, we have
replaced the original $m$
coordinate directions by
$k$ carefully selected
orthogonal directions.

Notice that $U_k$ contains
directions, not transformed
observations.

To obtain the new coordinates
of our observations, one
additional operation is required.

\subsubsection*{Step 4: Project the Data
onto the Principal Directions}

We compute

\[
\boxed{
Z=U_k^TA.
}
\]

Let us examine the dimensions:

\[
\underbrace{U_k^T}_{k\times m}
\underbrace{A}_{m\times n}
=
\underbrace{Z}_{k\times n}.
\]

The $j$th column of $Z$ is

\[
z_j=U_k^Ta_j.
\]

In particular,

\[
z_j=
\begin{bmatrix}
u_1^Ta_j\\
u_2^Ta_j\\
\vdots\\
u_k^Ta_j
\end{bmatrix}.
\]

Thus, each entry measures
how much the observation
extends along a particular
principal direction.

The original observation
$a_j\in\mathbb{R}^m$
has now been represented
by a new vector

\[
z_j\in\mathbb{R}^k.
\]

This is the actual
dimensionality reduction.

The matrix $Z$ is commonly
called the matrix of
\textit{principal component scores}.

An interesting observation
follows directly from SVD.
Since

\[
A=U\Sigma V^T,
\]

and the columns of $U$
are orthonormal,

\[
\boxed{
Z=U_k^TA=\Sigma_k V_k^T,
}
\]

where $\Sigma_k$ contains
the first $k$ positive
singular values and $V_k$
contains the corresponding
right singular vectors.

Therefore, the principal
component scores are precisely
the corresponding right
singular-vector coordinates
scaled by the singular values.

\subsubsection*{Step 5: Reconstruct the
Original Observations}

Suppose we now wish to
return from the reduced
$k$-dimensional representation
to the original feature space.

We compute

\[
\boxed{
\widehat{A}=U_kZ.
}
\]

Substituting $Z=U_k^TA$,

\[
\boxed{
\widehat{A}=U_kU_k^TA.
}
\]

Notice the dimensions:

\[
\underbrace{U_k}_{m\times k}
\underbrace{Z}_{k\times n}
=
\underbrace{\widehat{A}}_{m\times n}.
\]

Thus, the reconstructed
data matrix has the same
dimensions as the original
centered data matrix.

However, in general,

\[
\widehat{A}\neq A.
\]

The reconstructed observations
retain only the components
lying in the selected principal
subspace.

The components orthogonal to
that subspace are discarded.

If we wish to reconstruct
the observations on their
original, uncentered scale,
we must also restore the mean:

\[
\boxed{
\widehat{X}
=
\widehat{A}
+
\bar{x}\mathbf{1}_n^T.
}
\]

This completes the basic
PCA procedure.

% -------------------------------------------------
\subsection{What Have We Actually Accomplished?}
% -------------------------------------------------

It is worth reflecting on what
these operations have achieved.

Originally, every observation
was represented by $m$
features having familiar
interpretations, such as age,
income, and education.

After the transformation

\[
Z=U_k^TA,
\]

the observations are described
using only $k$ coordinates.

However, the new coordinates
no longer represent individual
features in the original sense.

For example, a principal
component might be a linear
combination of age, income,
and education, rather than
any one of these quantities
individually.

Thus, dimensionality reduction
comes with a trade-off:
the new coordinates may be
mathematically more economical,
but their interpretation
is often less immediate.

Nevertheless, the transformation
is not mysterious.

The coefficients defining
each new coordinate are
stored explicitly in the
columns of $U_k$.

When we reconstruct the
observations through

\[
\widehat{A}=U_kZ,
\]

we return to the original
feature coordinates.

The rows of $\widehat{A}$
can therefore again be
interpreted as age, income,
education, and so forth.

Their numerical values,
however, are generally
different from those
in the original matrix.

Geometrically, each observation
has been replaced by its
orthogonal projection onto
the principal subspace.

This process may suppress
small-scale fluctuations
and can be useful for
denoising when the discarded
directions predominantly
contain noise.

However, PCA does not
automatically distinguish
noise from meaningful
variation. Information
contained in discarded
directions is lost,
whether it represents
noise or an important
feature of the dataset.

% -------------------------------------------------
\subsection{Why Do We Retain the Largest
Singular Values?}
% -------------------------------------------------

We have now described the
PCA algorithm, but have
not yet established its
mathematical justification.

In particular, why should
we retain the left singular
vectors corresponding to
the largest singular values?

Let us begin with a single
unit vector

\[
u\in\mathbb{R}^m,
\qquad
u^Tu=1.
\]

Suppose we project every
centered observation onto
this direction.

The projected coordinate
of the $j$th observation is

\[
z_j=u^Ta_j.
\]

Since the original data
is centered, the projected
coordinates also have
mean zero.

Their sample variance is

\[
\operatorname{Var}(u^TA)
=
\frac{1}{n-1}
\sum_{j=1}^{n}(u^Ta_j)^2,
\qquad n\geq 2.
\]

Notice that

\[
\begin{aligned}
\sum_{j=1}^{n}(u^Ta_j)^2
&=
\sum_{j=1}^{n}
u^Ta_ja_j^Tu\\
&=
u^T
\left(
\sum_{j=1}^{n}a_ja_j^T
\right)u\\
&=u^TAA^Tu.
\end{aligned}
\]

Consequently,

\[
\boxed{
\operatorname{Var}(u^TA)
=
\frac{1}{n-1}u^TAA^Tu.
}
\]

We can therefore formulate
the problem of finding a
direction of maximum variance
as follows:

\[
\boxed{
\max_{\|u\|_2=1}
u^TAA^Tu.
}
\]

Now recall that the left
singular vectors of $A$
are eigenvectors of $AA^T$.

Indeed,

\[
AA^Tu_i=\sigma_i^2u_i.
\]

By the spectral theorem,
we may expand any unit
vector $u$ as

\[
u=\sum_{i=1}^{m}c_i u_i,
\qquad
\sum_{i=1}^{m}c_i^2=1,
\]

where zero eigenvalues are
included when necessary.

Therefore,

\[
u^TAA^Tu
=
\sum_{i=1}^{m}
\sigma_i^2c_i^2,
\]

where singular values beyond
the rank are understood
to be zero.

Since

\[
\sigma_1^2\geq\sigma_i^2,
\]

we immediately obtain

\[
\begin{aligned}
u^TAA^Tu
&\leq
\sigma_1^2
\sum_{i=1}^{m}c_i^2\\
&=\sigma_1^2.
\end{aligned}
\]

Equality is attained by
choosing $u=u_1$.

Hence,

\[
\boxed{
\max_{\|u\|_2=1}
u^TAA^Tu
=
\sigma_1^2.
}
\]

We have established an
important result:

\textit{A left singular vector
corresponding to the largest
singular value identifies
a direction of maximum
variance in the centered
dataset.}

The first principal direction
is therefore not chosen
arbitrarily.

It solves a well-defined
optimization problem.

The same reasoning can be
extended to selecting
several mutually orthogonal
directions.

The first $k$ left singular
vectors maximize the total
variance captured among all
choices of $k$ orthonormal
directions.

We will examine this
higher-dimensional statement
more carefully in the
subsequent mathematical
development.

% -------------------------------------------------
\subsection{The Connection with Reconstruction Error}
% -------------------------------------------------

There is another way to
understand the same result.

Suppose we choose any
unit vector $u$ and
reconstruct each observation
from its scalar projection
onto that direction.

For the $j$th observation,

\[
\widehat{a}_j
=
u(u^Ta_j).
\]

The reconstruction error is

\[
e_j=a_j-\widehat{a}_j.
\]

Because this is an
orthogonal projection,
the error vector is
perpendicular to the
reconstructed vector.

The Pythagorean theorem
therefore gives

\[
\|a_j\|_2^2
=
\|\widehat{a}_j\|_2^2
+
\|a_j-\widehat{a}_j\|_2^2.
\]

Since $\|u\|_2=1$,

\[
\|\widehat{a}_j\|_2^2
=
(u^Ta_j)^2.
\]

It follows that

\[
\|a_j-\widehat{a}_j\|_2^2
=
\|a_j\|_2^2-(u^Ta_j)^2.
\]

Summing over all observations
and using the Frobenius norm,

\[
\boxed{
\|A-uu^TA\|_F^2
=
\|A\|_F^2-u^TAA^Tu.
}
\]

Here,

\[
\|A\|_F^2
=
\sum_{i,j}A_{ij}^2.
\]

The quantity $\|A\|_F^2$
does not depend on $u$.

Consequently, minimizing
the reconstruction error
is equivalent to maximizing
the retained variance.

Thus,

\[
\boxed{
\begin{aligned}
&\max_{\|u\|_2=1}
u^TAA^Tu\\[2mm]
&\qquad\Longleftrightarrow\\[2mm]
&\min_{\|u\|_2=1}
\|A-uu^TA\|_F^2.
\end{aligned}
}
\]

This equivalence reveals
the geometric beauty
underlying PCA.

The direction that captures
the greatest variance is
also the direction that
produces the smallest
squared reconstruction error
among all one-dimensional
orthogonal projections.

More generally, if
$Q\in\mathbb{R}^{m\times k}$
has orthonormal columns,
then

\[
\boxed{
\|A-QQ^TA\|_F^2
=
\|A\|_F^2-\|Q^TA\|_F^2.
}
\]

Thus, maximizing the total
variation retained in the
$k$-dimensional subspace
is equivalent to minimizing
its squared reconstruction
error.

For PCA, we choose $Q=U_k$.

% -------------------------------------------------
\subsection{The Optimality Question}
% -------------------------------------------------

We have obtained the PCA
reconstruction

\[
\widehat{A}=U_kU_k^TA.
\]

Using the singular value
decomposition, this can
also be written as

\[
\boxed{
\widehat{A}
=
\sum_{i=1}^{k}
\sigma_i u_i v_i^T.
}
\]

Since the singular vectors
are orthonormal, the
discarded part is

\[
A-\widehat{A}
=
\sum_{i=k+1}^{r}
\sigma_i u_i v_i^T.
\]

Consequently, its squared
Frobenius norm is

\[
\boxed{
\|A-\widehat{A}\|_F^2
=
\sum_{i=k+1}^{r}
\sigma_i^2.
}
\]

Similarly, when $k<r$,
the spectral norm of the
residual is

\[
\boxed{
\|A-\widehat{A}\|_2
=
\sigma_{k+1}.
}
\]

These formulas tell us
exactly how much error
is introduced by discarding
the smaller singular values.

But an important question
still remains.

We have argued that PCA
selects an optimal
$k$-dimensional orthogonal
projection.

What if we do not restrict
ourselves to orthogonal
projections?

Suppose we allow any
matrix $B$ of rank at most
$k$ to approximate $A$.

Could some other construction
produce a smaller error?

Remarkably, the answer is no.

The Eckart--Young theorem
establishes that the truncated
SVD minimizes the approximation
error among all matrices
of rank at most $k$,
both in the Frobenius
and spectral norms.

In particular,

\[
\boxed{
\min_{\operatorname{rank}(B)\leq k}
\|A-B\|_F^2
=
\sum_{i=k+1}^{r}\sigma_i^2,
}
\]

and, for $k<r$,

\[
\boxed{
\min_{\operatorname{rank}(B)\leq k}
\|A-B\|_2
=
\sigma_{k+1}.
}
\]

These are substantially
stronger claims than
the elementary projection
argument alone establishes.

We shall therefore devote
the later sections of this
note to developing and
proving these results
rigorously.

For now, the essential
observation is that PCA
is not simply a procedure
for discarding coordinates.

It is a geometrically
motivated optimization
method whose optimality
can be characterized
precisely through the
singular values of the
data matrix.

This is the mathematical
connection that we seek
to understand in the
remainder of this note.

\end{comment}


% =================================================
\section{Geometric Motivation for Principal
Component Analysis}
% =================================================


The fundamental idea of PCA is that a collection of
high-dimensional observations may exhibit most
of its variation along only a few directions.
If we can identify these directions, we may be
able to describe the dataset using substantially
fewer coordinates without losing too much of
its original structure.

But what exactly do we mean by a direction
that captures the variation of a dataset?
And how can we determine which direction
is the most important?

We begin with a simple geometric interpretation.

% -------------------------------------------------
\subsection{Representing a Dataset Geometrically}
% -------------------------------------------------

Suppose we have $n$ observations, each described
by $m$ features. We represent the dataset by

\[
X =
\begin{bmatrix}
| & | & & |\\
x_1 & x_2 & \cdots & x_n\\
| & | & & |
\end{bmatrix}
\in \mathbb{R}^{m\times n},
\]

where each column $x_j\in\mathbb{R}^m$
represents one observation, and each row
corresponds to a particular feature.

For example, if our observations are individuals,
the rows might represent age, income, height,
or years of education. Each individual is
therefore represented by a point in an
$m$-dimensional Euclidean space.

When $m=2$, the observations form a collection
of points in the plane. When $m=3$, they form
a cloud of points in three-dimensional space.

For larger $m$, direct visualization becomes
difficult, but the underlying geometry
remains unchanged.

Now imagine that this cloud of points is
not distributed uniformly in every direction.
Instead, it may be elongated along some
particular direction.

Intuitively, that direction should contain
more information about how the observations
differ from one another than a direction
along which the points are tightly clustered.

This observation provides the initial
motivation for PCA.

The idea is to identify a new system of
orthogonal coordinate axes that reflects
the dominant directions of variation
in the dataset.

Rather than describing each observation
using all the original features, we may
then describe it using only a few
coordinates along these new axes.

But before proceeding, we must clarify
what we actually mean by the variation
of a dataset.

% -------------------------------------------------
\subsection{What Does Variance Mean Geometrically?}
% -------------------------------------------------

In elementary statistics, variance measures
how widely observations are distributed
around their mean.

For a collection of scalar observations
$t_1,t_2,\ldots,t_n$, the sample variance
is defined by

\[
s^2=\frac{1}{n-1}
\sum_{j=1}^{n}(t_j-\bar{t})^2,
\]

where

\[
\bar{t}=\frac{1}{n}
\sum_{j=1}^{n}t_j.
\]

A large variance means that the observations
are widely spread around their mean,
whereas a small variance means that they
are concentrated relatively close to it.

This interpretation is straightforward
when the observations are real numbers.

But what happens when each observation
is a vector in $\mathbb{R}^m$?

In that case, we must first specify
the direction along which the spread
is to be measured.

To understand this, imagine a two-dimensional
dataset containing information about
the income and expenditure of several
individuals.

Suppose individuals with higher income
also tend to have higher expenditure.

The corresponding data points may then
form an elongated cloud approximately
along a diagonal direction.

If we project these observations onto
that diagonal line, their projected
coordinates may be widely spread.

On the other hand, if we project the
same observations onto a direction
perpendicular to that line, their
projected coordinates may lie much
closer together.

Thus, the variance of the projected
observations depends on the direction
along which we examine the data.

This gives us a useful geometric
interpretation:

\textit{Variance along a direction
measures how widely the observations
are spread when viewed from that
particular direction.}

Let $u\in\mathbb{R}^m$ be an arbitrary
unit vector, so that

\[
u^Tu=1.
\]

For each observation $x_j$, consider
its scalar projection onto $u$:

\[
t_j=u^Tx_j.
\]

The mean of these projected coordinates
is

\[
\bar{t}
=
u^T\bar{x},
\]

where

\[
\bar{x}=\frac{1}{n}
\sum_{j=1}^{n}x_j.
\]

Therefore, the sample variance of the
observations along the direction $u$ is

\[
\boxed{
s_u^2=
\frac{1}{n-1}
\sum_{j=1}^{n}
\left(u^T(x_j-\bar{x})\right)^2.
}
\]

Notice something important.

Although the original dataset may
contain hundreds of features,
we can measure its spread along
any chosen direction by first
projecting the observations onto
that direction and then computing
the ordinary variance of the
resulting scalar coordinates.

This leads us naturally to a
mathematical question:

\textit{Among all possible unit
directions, which one produces
the greatest variance?}

PCA is designed to answer precisely
this question.

The first principal direction
captures the maximum possible
variance. Subsequent principal
directions capture the largest
remaining variance while being
orthogonal to the previously
selected directions.

This is the essential geometric
principle behind PCA.

Of course, a direction with large
variance is not necessarily
important for every conceivable
application. PCA uses variance
as a measure of the dominant
structure of a dataset.

The mathematical problem is
therefore not to preserve every
possible notion of information,
but to preserve as much
\textit{variance} as possible
using a prescribed number
of orthogonal directions.






% -------------------------------------------------
\subsection{The Covariance Matrix}
% -------------------------------------------------

Having introduced the notion of variance along
a particular direction, we now recall another
important concept from elementary statistics:
the covariance matrix.

While variance measures the spread of a single
variable, covariance measures how two variables
vary together.

For two variables $X_i$ and $X_j$, with
observations $x_{i\ell}$ and $x_{j\ell}$,
their sample covariance is defined by

\[
\operatorname{Cov}(X_i,X_j)
=
\frac{1}{n-1}
\sum_{\ell=1}^{n}
(x_{i\ell}-\bar{x}_i)
(x_{j\ell}-\bar{x}_j).
\]

Notice that when $i=j$, covariance reduces
to the ordinary sample variance:

\[
\operatorname{Cov}(X_i,X_i)
=
\operatorname{Var}(X_i).
\]

The covariance matrix collects all these
quantities into a single symmetric matrix:

\[
C=
\begin{bmatrix}
\operatorname{Var}(X_1) &
\cdots &
\operatorname{Cov}(X_1,X_m)\\
\vdots & \ddots & \vdots\\
\operatorname{Cov}(X_m,X_1) &
\cdots &
\operatorname{Var}(X_m)
\end{bmatrix}.
\]

Thus, the diagonal entries represent the
variances of individual features, while
the off-diagonal entries describe how
different features vary together.

Let us now see how this matrix arises
naturally from matrix multiplication.



\begin{comment}
    


\subsubsection*{A Simple $2\times 2$ Example}

Suppose we have two features and two
observations, represented by

\[
X=
\begin{bmatrix}
1 & 3\\
2 & 4
\end{bmatrix}.
\]

Following our convention, the columns
represent observations and the rows
represent features.

The mean of each feature is

\[
\bar{x}_1=\frac{1+3}{2}=2,
\qquad
\bar{x}_2=\frac{2+4}{2}=3.
\]

Subtracting these means from the
corresponding rows, we obtain the
centered data matrix

\[
A=
\begin{bmatrix}
1-2 & 3-2\\
2-3 & 4-3
\end{bmatrix}
=
\begin{bmatrix}
-1 & 1\\
-1 & 1
\end{bmatrix}.
\]

Now consider the matrix product

\[
AA^T=
\begin{bmatrix}
-1 & 1\\
-1 & 1
\end{bmatrix}
\begin{bmatrix}
-1 & -1\\
1 & 1
\end{bmatrix}.
\]

Performing the multiplication gives

\[
AA^T=
\begin{bmatrix}
2 & 2\\
2 & 2
\end{bmatrix}.
\]

What do these entries represent?

The first diagonal entry is the sum
of squared deviations of the first
feature. The second diagonal entry
has the same interpretation for the
second feature.

The off-diagonal entries are the sums
of products of deviations of the two
features from their respective means.

These are precisely the quantities
appearing in the definition of covariance.

Since $n=2$, the sample covariance
matrix is

\[
C=\frac{1}{n-1}AA^T
=
\begin{bmatrix}
2 & 2\\
2 & 2
\end{bmatrix}.
\]

The same reasoning applies to any
number of features and observations.

In general, for a centered data matrix
$A\in\mathbb{R}^{m\times n}$ with
$n\geq 2$, the sample covariance
matrix is

\[
\boxed{
C=\frac{1}{n-1}AA^T.
}
\]

Thus, the matrix $AA^T$, which we
encountered earlier in our discussion
of Singular Value Decomposition, is
closely connected to the statistical
variation of the dataset.

More precisely, it is the covariance
matrix multiplied by $n-1$.

This connection will be particularly
important in our development of PCA.

Indeed, for any unit direction
$u\in\mathbb{R}^m$, the variance of
the projected observations can be
written as

\[
\boxed{
\operatorname{Var}(u^TA)
=
u^TCu
=
\frac{1}{n-1}u^TAA^Tu.
}
\]

Consequently, finding directions
of maximum variance is closely
related to studying the eigenvectors
and eigenvalues of the covariance
matrix.

As we shall see, this is precisely
where Singular Value Decomposition
enters the mathematical development
of Principal Component Analysis.


\end{comment}










% -------------------------------------------------
\subsubsection*{A Simple $2\times 2$ Example}
% -------------------------------------------------

Suppose we have two features and two
observations, represented by the matrix

\[
X=
\begin{bmatrix}
x_{11} & x_{12}\\
x_{21} & x_{22}
\end{bmatrix}.
\]

Following our convention, the columns
represent observations and the rows
represent features. Thus, $x_{ij}$
denotes the value of the $i$th feature
in the $j$th observation.

The mean of each feature is

\[
\bar{x}_1=\frac{x_{11}+x_{12}}{2},
\qquad
\bar{x}_2=\frac{x_{21}+x_{22}}{2}.
\]

Subtracting the mean from each
corresponding row, we obtain
the centered data matrix

\[
A=
\begin{bmatrix}
x_{11}-\bar{x}_1 & x_{12}-\bar{x}_1\\
x_{21}-\bar{x}_2 & x_{22}-\bar{x}_2
\end{bmatrix}.
\]

For convenience, let us introduce
the notation

\[
a_{ij}=x_{ij}-\bar{x}_i.
\]

Then

\[
\boxed{
A=
\begin{bmatrix}
a_{11} & a_{12}\\
a_{21} & a_{22}
\end{bmatrix}.
}
\]

We will now calculate the covariance
matrix directly from the statistical
definitions, without using matrix
multiplication.

\subsubsection*{Computing the Covariance
Matrix}

Recall that the covariance matrix
for two features has the form

\[
C=
\begin{bmatrix}
\operatorname{Var}(X_1) &
\operatorname{Cov}(X_1,X_2)\\
\operatorname{Cov}(X_2,X_1) &
\operatorname{Var}(X_2)
\end{bmatrix}.
\]

Let us calculate each entry separately.

\textbf{First diagonal entry:}

The sample variance of the first
feature is

\[
\begin{aligned}
C_{11}
&=\operatorname{Var}(X_1)\\
&=\frac{1}{2-1}
\left[
(x_{11}-\bar{x}_1)^2
+
(x_{12}-\bar{x}_1)^2
\right]\\
&=a_{11}^2+a_{12}^2.
\end{aligned}
\]

\textbf{Second diagonal entry:}

Similarly, the sample variance
of the second feature is

\[
\begin{aligned}
C_{22}
&=\operatorname{Var}(X_2)\\
&=\frac{1}{2-1}
\left[
(x_{21}-\bar{x}_2)^2
+
(x_{22}-\bar{x}_2)^2
\right]\\
&=a_{21}^2+a_{22}^2.
\end{aligned}
\]

\textbf{First off-diagonal entry:}

The sample covariance between
the first and second features is

\[
\begin{aligned}
C_{12}
&=\operatorname{Cov}(X_1,X_2)\\
&=\frac{1}{2-1}
\Big[
(x_{11}-\bar{x}_1)
(x_{21}-\bar{x}_2)\\
&\qquad\qquad+
(x_{12}-\bar{x}_1)
(x_{22}-\bar{x}_2)
\Big]\\
&=a_{11}a_{21}+a_{12}a_{22}.
\end{aligned}
\]

\textbf{Second off-diagonal entry:}

By the symmetry of covariance,

\[
\begin{aligned}
C_{21}
&=\operatorname{Cov}(X_2,X_1)\\
&=a_{21}a_{11}+a_{22}a_{12}\\
&=C_{12}.
\end{aligned}
\]

Putting these four entries together,
we obtain the covariance matrix

\[
\boxed{
C=
\begin{bmatrix}
a_{11}^2+a_{12}^2 &
a_{11}a_{21}+a_{12}a_{22}\\[2mm]
a_{21}a_{11}+a_{22}a_{12} &
a_{21}^2+a_{22}^2
\end{bmatrix}.
}
\]

So far, we have used nothing
more than the elementary
statistical definitions of
variance and covariance.

Now comes an interesting observation.

% -------------------------------------------------
\subsubsection*{The Same Calculation Through
Matrix Multiplication}
% -------------------------------------------------

Recall that the centered
data matrix is

\[
A=
\begin{bmatrix}
a_{11} & a_{12}\\
a_{21} & a_{22}
\end{bmatrix}.
\]

Consider the matrix product

\[
AA^T=
\begin{bmatrix}
a_{11} & a_{12}\\
a_{21} & a_{22}
\end{bmatrix}
\begin{bmatrix}
a_{11} & a_{21}\\
a_{12} & a_{22}
\end{bmatrix}.
\]

Performing the multiplication,
we obtain

\[
\boxed{
AA^T=
\begin{bmatrix}
a_{11}^2+a_{12}^2 &
a_{11}a_{21}+a_{12}a_{22}\\[2mm]
a_{21}a_{11}+a_{22}a_{12} &
a_{21}^2+a_{22}^2
\end{bmatrix}.
}
\]

Compare this matrix with the
covariance matrix we calculated
earlier.

They are exactly the same!

In this particular example,
the number of observations
is $n=2$, so the sample
covariance normalization
factor is

\[
\frac{1}{n-1}=1.
\]

Hence,

\[
C=AA^T.
\]

But this equality without
a normalization factor is
specific to our example.

For a general dataset containing
$n$ observations, the sample
covariance formula requires
division by $n-1$.

Thus, for a centered data
matrix $A\in\mathbb{R}^{m\times n}$,
with $n\geq 2$, we obtain

\[
\boxed{
C=\frac{1}{n-1}AA^T.
}
\]

To see this directly, the
$(i,j)$ entry of $AA^T$ is

\[
(AA^T)_{ij}
=
\sum_{\ell=1}^{n}
a_{i\ell}a_{j\ell}.
\]

Therefore,

\[
\begin{aligned}
C_{ij}
&=\frac{1}{n-1}(AA^T)_{ij}\\
&=\frac{1}{n-1}
\sum_{\ell=1}^{n}
(x_{i\ell}-\bar{x}_i)
(x_{j\ell}-\bar{x}_j).
\end{aligned}
\]

This is precisely the
definition of sample covariance
between the $i$th and $j$th
features.

We have therefore established
the covariance matrix formula
directly from its elementary
statistical definition.

% -------------------------------------------------
\subsubsection*{Connection with PCA}
% -------------------------------------------------

The matrix $AA^T$, which we
encountered earlier in our
discussion of Singular Value
Decomposition, is therefore
closely connected to the
statistical variation of
the dataset.

More precisely, it is the
covariance matrix multiplied
by $n-1$.

This connection is especially
important for PCA.


Let $u\in\mathbb{R}^m$ be a unit vector,
and let the centered data matrix be

\[
A=
\begin{bmatrix}
a_1 & a_2 & \cdots & a_n
\end{bmatrix}.
\]

The scalar projection of the $j$th
observation onto the direction $u$ is

\[
z_j=u^Ta_j.
\]

Since the observations have been centered,
the mean of the projected coordinates is

\[
\begin{aligned}
\bar{z}
&=\frac{1}{n}\sum_{j=1}^{n}z_j\\
&=\frac{1}{n}\sum_{j=1}^{n}u^Ta_j\\
&=u^T\left(\frac{1}{n}
\sum_{j=1}^{n}a_j\right)\\
&=0.
\end{aligned}
\]

Therefore, by the definition of
sample variance,

\[
\begin{aligned}
\operatorname{Var}(u^TA)
&=\frac{1}{n-1}
\sum_{j=1}^{n}(z_j-\bar{z})^2\\[2mm]
&=\frac{1}{n-1}
\sum_{j=1}^{n}(u^Ta_j)^2\\[2mm]
&=\frac{1}{n-1}
\sum_{j=1}^{n}
(u^Ta_j)(a_j^Tu)\\[2mm]
&=\frac{1}{n-1}
\sum_{j=1}^{n}
u^T(a_ja_j^T)u\\[2mm]
&=\frac{1}{n-1}
u^T\left(
\sum_{j=1}^{n}a_ja_j^T
\right)u\\[2mm]
&=\frac{1}{n-1}u^TAA^Tu\\[2mm]
&=u^TCu,
\end{aligned}
\]

where we have used the identity

\[
AA^T=\sum_{j=1}^{n}a_ja_j^T
\]

and the covariance matrix formula

\[
C=\frac{1}{n-1}AA^T.
\]

Hence, we have established

\[
\boxed{
\operatorname{Var}(u^TA)
=
u^TCu
=
\frac{1}{n-1}u^TAA^Tu.
}
\]

This is a fundamental relationship
in Principal Component Analysis.

It tells us that the variance of
the observations along an arbitrary
unit direction $u$ can be expressed
as a quadratic form involving
the covariance matrix.

Consequently, finding the direction
of maximum variance is equivalent
to solving the optimization problem

\[
\boxed{
\max_{u^Tu=1}u^TCu.
}
\]

As we shall establish shortly,
the solution is an eigenvector
of $C$ corresponding to its
largest eigenvalue.




Thus, we can measure how much
the data varies along any
particular direction using
a quadratic form involving
the covariance matrix.



As we shall see, this leads
naturally to the eigenvectors
of the covariance matrix,
or equivalently, the left
singular vectors of the
centered data matrix.

This is precisely where
the statistical concept
of variance meets the
linear algebra of Singular
Value Decomposition.


% -------------------------------------------------
\subsection{The Principal Component Analysis
Procedure}
% -------------------------------------------------

As we have developed the necessary mathematical tools, the basic PCA procedure can now be summarized
in a few relatively simple steps.

Although different textbooks may arrange
the data matrix differently, we will
consistently assume that features occupy
the rows and observations occupy the
columns.

\subsubsection*{Step 1: Center the Data}

We begin by computing the mean
observation vector

\[
\bar{x}=
\frac{1}{n}\sum_{j=1}^{n}x_j.
\]

Each observation is then centered
by subtracting this mean:

\[
a_j=x_j-\bar{x}.
\]

The resulting centered data matrix is

\[
\boxed{
A=X-\bar{x}\mathbf{1}_n^T,
}
\]

where $\mathbf{1}_n$ is the vector
of $n$ ones.

Notice that

\[
\sum_{j=1}^{n}a_j=0.
\]

Consequently, the center of the
data cloud has been translated
to the origin.

Why do we perform this operation?

Recall that variance measures
the spread of observations
\textit{around their mean}.

If we wish to identify directions
of maximum variance, we should
therefore examine the deviations
from the mean rather than the
absolute positions of the
observations.

Mean centering allows us to
measure this variation directly.

Indeed, for any unit vector $u$,
the variance along $u$ becomes

\[
s_u^2=
\frac{1}{n-1}
\sum_{j=1}^{n}(u^Ta_j)^2.
\]

We have already observed that this
expression has a particularly
convenient matrix representation.

From now on, unless stated
otherwise, $A$ denotes the
centered data matrix.

\subsubsection*{Step 2: Compute the SVD}

We compute the Singular Value
Decomposition of $A$:

\[
\boxed{
A=U\Sigma V^T.
}
\]

Recall from Section 3 that
the columns of $U$ are
the left singular vectors,
while the columns of $V$
are the right singular vectors.

The positive singular values
are arranged in decreasing order:

\[
\sigma_1\geq\sigma_2
\geq\cdots\geq\sigma_r>0,
\]

where $r=\operatorname{rank}(A)$.

Since $A$ is an $m\times n$
matrix, its left singular
vectors belong to $\mathbb{R}^m$,
the original feature space.

Thus, these vectors provide
candidate directions along
which we may describe the
observations.

But why should SVD help us
identify directions of
maximum variance?

To appreciate this connection,
recall the sample covariance
matrix of the centered data:

\[
\boxed{
C=\frac{1}{n-1}AA^T,
}
\]

where we assume $n\geq 2$.

The covariance matrix describes
how the original features vary
together.

Its diagonal entries are
the sample variances of the
individual features, while
the off-diagonal entries
measure their pairwise
covariances.

More importantly, for any
unit direction $u$,

\[
\boxed{
s_u^2=u^TCu.
}
\]

We have already derived this identity
explicitly in the previous subsection.

Now recall from SVD that

\[
AA^T
=
U\Sigma\Sigma^TU^T.
\]

Consequently,

\[
C=
\frac{1}{n-1}
U\Sigma\Sigma^TU^T.
\]

Thus, the left singular vectors
of $A$ are eigenvectors of
the covariance matrix $C$,
with corresponding eigenvalues

\[
\boxed{
\lambda_i=
\frac{\sigma_i^2}{n-1}.
}
\]

This is the first important
connection between SVD and
variance.

The squared singular values,
after division by $n-1$,
are precisely the variances
associated with the corresponding
principal directions.

Therefore, larger singular
values correspond to directions
along which the projected
observations exhibit greater
variance.

We will establish this
conclusion rigorously below.

\subsubsection*{Step 3: Select the Principal
Directions}

Suppose we wish to represent
the original $m$-dimensional
observations using only
$k$ coordinates, where

\[
1\leq k\leq r.
\]

We retain the first $k$
columns of $U$:

\[
\boxed{
U_k=
\begin{bmatrix}
u_1 & u_2 & \cdots & u_k
\end{bmatrix}
\in\mathbb{R}^{m\times k}.
}
\]

These vectors form an
orthonormal basis for a
$k$-dimensional subspace
of the original feature space.

We call this the
\textit{principal subspace}.

Why do we retain the first
$k$ vectors rather than
some other $k$ vectors?

Because the singular values
are arranged in decreasing order,
the selected vectors correspond
to the largest variances

\[
\frac{\sigma_1^2}{n-1},
\ldots,
\frac{\sigma_k^2}{n-1}.
\]

Thus, PCA retains those
orthogonal directions along
which the observations exhibit
the greatest total variation.




\begin{comment}
The total sample variance
across all original features is

\[
\operatorname{tr}(C)
=
\frac{1}{n-1}\|A\|_F^2
=
\frac{1}{n-1}
\sum_{i=1}^{r}\sigma_i^2.
\]

\end{comment}


The total sample variance
across all original features is

\[
\begin{aligned}
\operatorname{tr}(C)
&=\operatorname{tr}
\left(\frac{1}{n-1}AA^T\right)\\[2mm]
&=\frac{1}{n-1}\operatorname{tr}(AA^T)\\[2mm]
&=\frac{1}{n-1}
\operatorname{tr}
\left(U\Sigma V^T V\Sigma^TU^T\right)\\[2mm]
&=\frac{1}{n-1}
\operatorname{tr}
\left(U\Sigma\Sigma^TU^T\right)\\[2mm]
&=\frac{1}{n-1}
\operatorname{tr}
\left(\Sigma\Sigma^TU^TU\right)\\[2mm]
&=\frac{1}{n-1}
\operatorname{tr}\left(\Sigma\Sigma^T\right)\\[2mm]
&=\frac{1}{n-1}
\sum_{i=1}^{r}\sigma_i^2\\[2mm]
&=\frac{1}{n-1}\|A\|_F^2.
\end{aligned}
\]

Here we have used the orthogonality
of $U$ and $V$, together with the
cyclic property of the trace,

\[
\operatorname{tr}(BC)
=
\operatorname{tr}(CB),
\]

whenever both products are defined.

We have also used the identity

\[
\|A\|_F^2
=
\operatorname{tr}(AA^T)
=
\sum_{i=1}^{r}\sigma_i^2.
\]

Consequently,

\[
\boxed{
\operatorname{tr}(C)
=
\frac{1}{n-1}\|A\|_F^2
=
\frac{1}{n-1}
\sum_{i=1}^{r}\sigma_i^2.
}
\]


For a nonzero centered dataset,
the proportion of this variance
retained by the first $k$
principal directions is

\[
\boxed{
\mathrm{EVR}_k
=
\frac{\displaystyle
\sum_{i=1}^{k}\sigma_i^2}
{\displaystyle
\sum_{i=1}^{r}\sigma_i^2}.
}
\]

This quantity is called the
\textit{cumulative explained
variance ratio}.

It allows us to measure
how much of the original
variation is retained
after dimensionality
reduction.

Notice that $U_k$ contains
directions, not transformed
observations.

To obtain the new coordinates
of our observations, one
additional operation is required.

\subsubsection*{Step 4: Project the Data
onto the Principal Directions}

We compute

\[
\boxed{
Z=U_k^TA.
}
\]

Let us examine the dimensions:

\[
\underbrace{U_k^T}_{k\times m}
\underbrace{A}_{m\times n}
=
\underbrace{Z}_{k\times n}.
\]

The $j$th column of $Z$ is

\[
z_j=U_k^Ta_j.
\]

In particular,

\[
z_j=
\begin{bmatrix}
u_1^Ta_j\\
u_2^Ta_j\\
\vdots\\
u_k^Ta_j
\end{bmatrix}.
\]

Thus, each entry measures
how much the observation
extends along a particular
principal direction.

The original observation
$a_j\in\mathbb{R}^m$
has now been represented
by a new vector

\[
z_j\in\mathbb{R}^k.
\]

This is the actual
dimensionality reduction.

The matrix $Z$ is commonly
called the matrix of
\textit{principal component scores}.

An interesting observation
follows directly from SVD.
Since

\[
A=U\Sigma V^T,
\]

and the columns of $U$
are orthonormal,

\[
\boxed{
Z=U_k^TA=\Sigma_kV_k^T,
}
\]

where $\Sigma_k$ is the
$k\times k$ diagonal matrix
of the retained singular values,
and $V_k$ contains the corresponding
right singular vectors.

Therefore, the principal
component scores are precisely
the corresponding right
singular-vector coordinates
scaled by the singular values.

Now let us return once more
to the concept of variance.

The $i$th row of $Z$ contains
the projected coordinates of
all observations along
the direction $u_i$.

Since the data is centered,
the sample variance of
these coordinates is

\[
\operatorname{Var}(\mathrm{PC}_i)
=
\frac{1}{n-1}
\|u_i^TA\|_2^2.
\]

But

\[
u_i^TA=\sigma_i v_i^T.
\]

Since $\|v_i\|_2=1$,

\[
\boxed{
\operatorname{Var}(\mathrm{PC}_i)
=
\frac{\sigma_i^2}{n-1}.
}
\]

Thus, the singular values
tell us exactly how much
variance is captured
by each principal component.

There is another important
property of the transformed data.

Since

\[
Z=\Sigma_kV_k^T,
\]

we obtain

\[
\frac{1}{n-1}ZZ^T
=
\frac{1}{n-1}\Sigma_k^2.
\]

Hence, the covariance matrix
of the principal component
scores is diagonal:

\[
\boxed{
\operatorname{Cov}(Z)
=
\frac{1}{n-1}
\begin{bmatrix}
\sigma_1^2 & & 0\\
&\ddots&\\
0&&\sigma_k^2
\end{bmatrix}.
}
\]

Consequently, the principal
components are mutually
uncorrelated.

This gives another interpretation
of PCA: it constructs new
uncorrelated coordinates,
ordered according to the
amount of variance they capture.

\subsubsection*{Step 5: Reconstruct the
Original Observations}

Suppose we now wish to
return from the reduced
$k$-dimensional representation
to the original feature space.

We compute

\[
\boxed{
\widehat{A}=U_kZ.
}
\]

Substituting $Z=U_k^TA$,

\[
\boxed{
\widehat{A}=U_kU_k^TA.
}
\]

Notice the dimensions:

\[
\underbrace{U_k}_{m\times k}
\underbrace{Z}_{k\times n}
=
\underbrace{\widehat{A}}_{m\times n}.
\]

Thus, the reconstructed
data matrix has the same
dimensions as the original
centered data matrix.

However, in general,

\[
\widehat{A}\neq A.
\]

The reconstructed observations
retain only the components
lying in the selected principal
subspace.

The components orthogonal to
that subspace are discarded.

If we wish to reconstruct
the observations on their
original, uncentered scale,
we must also restore the mean:

\[
\boxed{
\widehat{X}
=
\widehat{A}
+
\bar{x}\mathbf{1}_n^T.
}
\]

This completes the basic
PCA procedure.


    

% -------------------------------------------------
\subsection{What Have We Actually Accomplished?}
% -------------------------------------------------

It is worth reflecting on what
these operations have achieved.

Originally, every observation
was represented by $m$
features having familiar
interpretations, such as age,
income, and education.

After the transformation

\[
Z=U_k^TA,
\]

the observations are described
using only $k$ coordinates.

However, the new coordinates
no longer represent individual
features in the original sense.

For example, a principal
component might be a linear
combination of age, income,
and education, rather than
any one of these quantities
individually.

Thus, dimensionality reduction
comes with a trade-off:
the new coordinates may be
mathematically more economical,
but their interpretation
is often less immediate.

Nevertheless, the transformation
is not mysterious.

The coefficients defining
each new coordinate are
stored explicitly in the
columns of $U_k$.

When we reconstruct the
observations through

\[
\widehat{A}=U_kZ,
\]

we return to the original
feature coordinates.

The rows of $\widehat{A}$
can therefore again be
interpreted as age, income,
education, and so forth.

Their numerical values,
however, are generally
different from those
in the original matrix.

Geometrically, each observation
has been replaced by its
orthogonal projection onto
the principal subspace.

This process may suppress
small-scale fluctuations
and can be useful for
denoising when the discarded
directions predominantly
contain noise.

However, PCA does not
automatically distinguish
noise from meaningful
variation. Information
contained in discarded
directions is lost,
whether it represents
noise or an important
feature of the dataset.

We should also remember
that variance depends
on the scale of the features.

For example, features measured
in dollars may numerically
dominate features measured
in years. Therefore, depending
on the application, it may
be appropriate to standardize
the features before applying PCA.

Such standardization changes
the mathematical problem,
since PCA then operates on
the scaled data rather than
the original measurements.


% -------------------------------------------------
\subsection{The Connection with Reconstruction Error}
% -------------------------------------------------

There is another way to
understand the same result.

Suppose we choose any
unit vector $u$ and
reconstruct each observation
from its scalar projection
onto that direction.

For the $j$th observation,

\[
\widehat{a}_j
=
u(u^Ta_j).
\]

The reconstruction error is

\[
e_j=a_j-\widehat{a}_j.
\]

Because this is an
orthogonal projection,
the error vector is
perpendicular to the
reconstructed vector.

The Pythagorean theorem
therefore gives

\[
\|a_j\|_2^2
=
\|\widehat{a}_j\|_2^2
+
\|a_j-\widehat{a}_j\|_2^2.
\]

Since $\|u\|_2=1$,

\[
\|\widehat{a}_j\|_2^2
=
(u^Ta_j)^2.
\]

It follows that

\[
\|a_j-\widehat{a}_j\|_2^2
=
\|a_j\|_2^2-(u^Ta_j)^2.
\]

Summing over all observations
and using the Frobenius norm,

\[
\boxed{
\|A-uu^TA\|_F^2
=
\|A\|_F^2-u^TAA^Tu.
}
\]

Here,

\[
\|A\|_F^2
=
\sum_{i,j}A_{ij}^2.
\]

The quantity $\|A\|_F^2$
does not depend on $u$.

Consequently, minimizing
the reconstruction error
is equivalent to maximizing
the retained variance.

Thus,

\[
\boxed{
\begin{aligned}
&\max_{\|u\|_2=1}
u^TAA^Tu\\[2mm]
&\qquad\Longleftrightarrow\\[2mm]
&\min_{\|u\|_2=1}
\|A-uu^TA\|_F^2.
\end{aligned}
}
\]

This equivalence reveals
the geometric beauty
underlying PCA.

The direction that captures
the greatest variance is
also the direction that
produces the smallest
squared reconstruction error
among all one-dimensional
orthogonal projections.

More generally, if
$Q\in\mathbb{R}^{m\times k}$
has orthonormal columns,
then

\[
\boxed{
\|A-QQ^TA\|_F^2
=
\|A\|_F^2-\|Q^TA\|_F^2.
}
\]

Thus, maximizing the total
variation retained in the
$k$-dimensional subspace
is equivalent to minimizing
its squared reconstruction
error.

For PCA, we choose $Q=U_k$.

% -------------------------------------------------
\subsection{The Optimality Question}
% -------------------------------------------------

We have obtained the PCA
reconstruction

\[
\widehat{A}=U_kU_k^TA.
\]

Using the singular value
decomposition, this can
also be written as

\[
\boxed{
\widehat{A}
=
\sum_{i=1}^{k}
\sigma_i u_i v_i^T.
}
\]

Since the singular vectors
are orthonormal, the
discarded part is

\[
A-\widehat{A}
=
\sum_{i=k+1}^{r}
\sigma_i u_i v_i^T.
\]

Consequently, its squared
Frobenius norm is

\[
\boxed{
\|A-\widehat{A}\|_F^2
=
\sum_{i=k+1}^{r}
\sigma_i^2.
}
\]

Similarly, when $k<r$,
the spectral norm of the
residual is

\[
\boxed{
\|A-\widehat{A}\|_2
=
\sigma_{k+1}.
}
\]

These formulas tell us
exactly how much error
is introduced by discarding
the smaller singular values.

But an important question
still remains.

We have shown that PCA
selects an optimal
$k$-dimensional orthogonal
projection under the squared
Frobenius reconstruction error.

What if we do not restrict
ourselves to orthogonal
projections?

Suppose we allow any
matrix $B$ of rank at most
$k$ to approximate $A$.

Could some other construction
produce a smaller error?

Remarkably, the answer is no.

The Eckart--Young theorem
establishes that the truncated
SVD minimizes the approximation
error among all matrices
of rank at most $k$,
both in the Frobenius
and spectral norms.

In particular,

\[
\boxed{
\min_{\operatorname{rank}(B)\leq k}
\|A-B\|_F^2
=
\sum_{i=k+1}^{r}\sigma_i^2,
}
\]

and, for $k<r$,

\[
\boxed{
\min_{\operatorname{rank}(B)\leq k}
\|A-B\|_2
=
\sigma_{k+1}.
}
\]

These are substantially
stronger claims than
the elementary projection
argument alone establishes.

We shall therefore devote
the later sections of this
note to developing and
proving these results
rigorously.

For now, the essential
observation is that PCA
is not simply a procedure
for discarding coordinates.

It is a geometrically
motivated optimization
method whose optimality
can be characterized
precisely through the
singular values of the
data matrix.

This is the mathematical
connection that we seek
to understand in the
remainder of this note.


% =================================================
\section{Mathematical Justification of Optimal
Low-Rank Approximation: The Eckart--Young Theorem}
% =================================================

In the preceding section, we saw that PCA
selects a low-dimensional subspace and
reconstructs the centered data through
an orthogonal projection. We also found
that the squared reconstruction error is
determined by the singular values that
have been discarded.

However, there is an important distinction
between calculating the error of a particular
approximation and proving that no other
approximation of same rank can do better.

So far, we have focused on approximations
arising from orthogonal projections. What
if we consider every possible matrix of
rank at most $k$, without requiring it
to arise from a projection?

The Eckart--Young theorem answers this
stronger question. We will first state
the result and then prove it for two
important matrix norms.

% -------------------------------------------------
\subsection{The Approximation Problem and Matrix Norms}
% -------------------------------------------------

Let $A\in\mathbb{R}^{m\times n}$ be any
real matrix of rank $r$, with singular
value decomposition

\[
A=\sum_{i=1}^{r}\sigma_i u_i v_i^T,
\qquad
\sigma_1\geq\cdots\geq\sigma_r>0.
\]

For $0\leq k<r$, define the truncated
singular value decomposition by

\[
\boxed{
A_k=\sum_{i=1}^{k}\sigma_i u_i v_i^T,
}
\]

where the empty sum for $k=0$ is understood
to be the zero matrix. Because the singular
vectors are orthonormal, $A_k$ has rank
at most $k$ (in fact, rank $k$ when $k>0$).

Our question is whether there exists
some other matrix $B$ satisfying

\[
\operatorname{rank}(B)\leq k
\]

that approximates $A$ more accurately.
Of course, to make the word ``accurately''
mathematically meaningful, we must first
specify how approximation error is measured.

\textbf{The Frobenius norm.}

For $M=(m_{ij})$, define

\[
\boxed{
\|M\|_F
=\left(\sum_{i,j}m_{ij}^{2}\right)^{1/2}.
}
\]

This measures the total entrywise error.
In particular, $\|A-B\|_F^2$ is the sum
of the squared differences between all
corresponding entries of $A$ and $B$.

\textbf{The spectral norm.}

The matrix spectral norm, also called
the induced Euclidean operator norm, is

\[
\boxed{
\|M\|_2
=\max_{\|x\|_2=1}\|Mx\|_2.
}
\]

It measures the greatest length that
$M$ can produce when acting on a unit
vector. For an approximation error
$M=A-B$, it measures the largest
possible directional error.

Thus, these norms ask somewhat different
questions: the Frobenius norm measures
the total squared error across the matrix,
whereas the spectral norm measures the
worst error over all unit input directions.

% -------------------------------------------------
\subsection{Statement of the Eckart--Young Theorem}
% -------------------------------------------------

\begin{theorem}[Eckart--Young--Mirsky]

Let $A\in\mathbb{R}^{m\times n}$ have rank
$r$ and positive singular values
$\sigma_1\geq\cdots\geq\sigma_r>0$.
For an integer $0\leq k<r$, define

\[
A_k=\sum_{i=1}^{k}\sigma_i u_i v_i^T.
\]

Then, among all matrices $B$ of rank
at most $k$, the truncated SVD $A_k$
minimizes both the Frobenius norm
and the spectral norm of the error.
More precisely,

\[
\boxed{
\min_{\operatorname{rank}(B)\leq k}
\|A-B\|_F^2
=\sum_{i=k+1}^{r}\sigma_i^2,
}
\]

and

\[
\boxed{
\min_{\operatorname{rank}(B)\leq k}
\|A-B\|_2=\sigma_{k+1}.
}
\]

Both minima are attained by $B=A_k$.

\end{theorem}

The key phrase is \textit{among all
matrices of rank at most $k$}.

The theorem does not assume that $B$
is formed by projection, nor does it
require that $B$ share the singular
vectors of $A$.

To establish the theorem, we must show
two things for each norm: every admissible
$B$ has error at least as large as the
claimed bound, and the truncated SVD
actually attains that bound.

% -------------------------------------------------
\subsection{Proof Under the Frobenius Norm}
% -------------------------------------------------

\begin{proof}

The case $k=0$ is immediate, since the
only matrix of rank zero is $B=0$.
We therefore assume $1\leq k<r$.

\textbf{Step 1: Begin with an arbitrary
rank-$k$ approximation.}

Let $B\in\mathbb{R}^{m\times n}$ be
\emph{any} matrix with

\[
\operatorname{rank}(B)\leq k.
\]

Its column space has dimension at
most $k$. Since $k<m$, we can choose
an orthonormal matrix

\[
Q\in\mathbb{R}^{m\times k},
\qquad Q^TQ=I_k,
\]

such that the column space of $B$
is contained in $\operatorname{Col}(Q)$.

Set

\[
P=QQ^T.
\]

Then $P$ is the orthogonal projector
onto $\operatorname{Col}(Q)$, and
because every column of $B$ belongs
to this subspace,

\[
\boxed{PB=B.}
\]

\textbf{Step 2: Separate the approximation
error into orthogonal parts.}

Write

\[
\begin{aligned}
A-B
&=(A-PA)+(PA-B)\\
&=(I-P)A+(PA-B).
\end{aligned}
\]

Every column of $(I-P)A$ is
orthogonal to $\operatorname{Col}(Q)$,
whereas every column of $PA-B$
lies inside that subspace.

Consequently, these two matrices
are orthogonal in the Frobenius
inner product, and the Pythagorean
theorem gives

\[
\boxed{
\|A-B\|_F^2
=\|(I-P)A\|_F^2
+\|PA-B\|_F^2.
}
\]

In particular,

\[
\boxed{
\|A-B\|_F^2\geq\|(I-P)A\|_F^2.
}
\]

Thus, no matrix $B$ whose columns
lie in this subspace can approximate
$A$ more accurately than the
orthogonal projection $PA$.

\textbf{Step 3: Bound the variance that
an arbitrary $k$-dimensional subspace
can retain.}

Using the same orthogonal decomposition,

\[
\begin{aligned}
\|(I-P)A\|_F^2
&=\|A\|_F^2-\|PA\|_F^2\\
&=\|A\|_F^2-\|Q^TA\|_F^2.
\end{aligned}
\]

We now show that

\[
\boxed{
\|Q^TA\|_F^2\leq
\sum_{i=1}^{k}\sigma_i^2.
}
\]

Recall from SVD that

\[
AA^T=\sum_{i=1}^{r}\sigma_i^2u_i u_i^T.
\]

Therefore,

\[
\begin{aligned}
\|Q^TA\|_F^2
&=\operatorname{tr}(Q^TAA^TQ)\\
&=\sum_{i=1}^{r}\sigma_i^2
\operatorname{tr}(Q^Tu_i u_i^TQ)\\
&=\sum_{i=1}^{r}\sigma_i^2\|Q^Tu_i\|_2^2.
\end{aligned}
\]

Put $p_i=\|Q^Tu_i\|_2^2$ and extend
the singular values by $\sigma_i=0$
for $r<i\leq m$. Since $Q$ has
orthonormal columns,

\[
0\leq p_i\leq 1,
\qquad
\sum_{i=1}^{m}p_i=k.
\]

The last identity follows from

\[
\begin{aligned}
\sum_{i=1}^{m}p_i
&=\operatorname{tr}\left(
Q^T\left(\sum_{i=1}^{m}u_i u_i^T\right)Q
\right)\\
&=\operatorname{tr}(Q^TQ)=k.
\end{aligned}
\]




We now wish to establish the inequality

\[
\sum_{i=1}^{m}\sigma_i^2p_i
\leq
\sum_{i=1}^{k}\sigma_i^2.
\]



Let us divide the sum into two parts:

\[
\begin{aligned}
\sum_{i=1}^{m}\sigma_i^2p_i
&=
\sum_{i=1}^{k}\sigma_i^2p_i
+
\sum_{i=k+1}^{m}\sigma_i^2p_i.
\end{aligned}
\]

Since $\sigma_i^2\leq\sigma_k^2$
whenever $i>k$, and $p_i\geq 0$,
we have

\[
\sum_{i=k+1}^{m}\sigma_i^2p_i
\leq
\sigma_k^2\sum_{i=k+1}^{m}p_i.
\]

Consequently,

\[
\begin{aligned}
\sum_{i=1}^{m}\sigma_i^2p_i
&\leq
\sum_{i=1}^{k}\sigma_i^2p_i
+
\sigma_k^2\sum_{i=k+1}^{m}p_i.
\end{aligned}
\]

Now recall that

\[
\sum_{i=1}^{m}p_i=k.
\]

Therefore,

\[
\begin{aligned}
\sum_{i=k+1}^{m}p_i
&=k-\sum_{i=1}^{k}p_i\\
&=\sum_{i=1}^{k}(1-p_i).
\end{aligned}
\]

Substituting this expression
into the previous inequality gives

\[
\begin{aligned}
\sum_{i=1}^{m}\sigma_i^2p_i
&\leq
\sum_{i=1}^{k}\sigma_i^2p_i
+
\sigma_k^2\sum_{i=1}^{k}(1-p_i)\\
&=
\sum_{i=1}^{k}
\left[
\sigma_i^2p_i+
\sigma_k^2(1-p_i)
\right].
\end{aligned}
\]

Now here comes the crucial observation.

For every $1\leq i\leq k$,

\[
\sigma_i^2\geq\sigma_k^2.
\]

Also, since $0\leq p_i\leq 1$,

\[
1-p_i\geq 0.
\]

Hence,

\[
\sigma_k^2(1-p_i)
\leq
\sigma_i^2(1-p_i).
\]

Adding $\sigma_i^2p_i$
to both sides, we obtain

\[
\begin{aligned}
\sigma_i^2p_i+\sigma_k^2(1-p_i)
&\leq
\sigma_i^2p_i+\sigma_i^2(1-p_i)\\
&=
\sigma_i^2[p_i+(1-p_i)]\\
&=\sigma_i^2.
\end{aligned}
\]

Summing this inequality
from $i=1$ to $k$,
we finally obtain

\[
\boxed{
\sum_{i=1}^{m}\sigma_i^2p_i
\leq
\sum_{i=1}^{k}\sigma_i^2.
}
\]

Equality is attained by choosing
$Q=U_k$, for which

\[
p_i=
\begin{cases}
1, & 1\leq i\leq k,\\
0, & i>k.
\end{cases}
\]

Thus, the first $k$ left singular
vectors maximize the quantity

\[
\|Q^TA\|_F^2
=
\sum_{i=1}^{m}\sigma_i^2p_i
\]

over all matrices $Q$ with
orthonormal columns.

This proves the claimed bound on
$\|Q^TA\|_F^2$.

\textbf{Step 4: Obtain a lower bound
for every admissible $B$.}

Combining our inequalities,

\[
\begin{aligned}
\|A-B\|_F^2
&\geq \|(I-P)A\|_F^2\\
&=\|A\|_F^2-\|Q^TA\|_F^2\\
&\geq\sum_{i=1}^{r}\sigma_i^2
-\sum_{i=1}^{k}\sigma_i^2\\
&=\sum_{i=k+1}^{r}\sigma_i^2.
\end{aligned}
\]

This inequality holds for \emph{every}
$B$ satisfying $\operatorname{rank}(B)\leq k$.

\textbf{Step 5: Show that the truncated
SVD achieves the bound.}

For

\[
A_k=\sum_{i=1}^{k}\sigma_i u_i v_i^T,
\]

we have

\[
A-A_k=\sum_{i=k+1}^{r}\sigma_i u_i v_i^T.
\]

The rank-one matrices $u_i v_i^T$
are orthonormal under the Frobenius
inner product. Therefore,

\[
\boxed{
\|A-A_k\|_F^2
=\sum_{i=k+1}^{r}\sigma_i^2.
}
\]

The lower bound is attained, proving

\[
\boxed{
\min_{\operatorname{rank}(B)\leq k}
\|A-B\|_F^2
=\sum_{i=k+1}^{r}\sigma_i^2.
}
\]

\end{proof}

Notice what made this proof work.

We did not begin by assuming that
$B$ had the same singular vectors
as $A$. Instead, we allowed an
arbitrary rank-$k$ matrix and showed
that an orthogonal projection onto
its column space, or a larger
$k$-dimensional subspace, could
never have a greater error than $B$.

We then used SVD to bound how much
variation any such subspace can retain.

% -------------------------------------------------
\subsection{Proof Under the Spectral Norm}
% -------------------------------------------------

The spectral-norm proof uses a different
idea. Rather than summing errors over
all entries, we need only identify
\emph{one} unit vector on which every
rank-$k$ approximation must incur an
error of at least $\sigma_{k+1}$.

\begin{proof}

Let $B\in\mathbb{R}^{m\times n}$ be
any matrix with $\operatorname{rank}(B)\leq k$.

\textbf{Step 1: Choose a subspace of
dimension $k+1$.}

Consider the first $k+1$ right
singular vectors of $A$ and define

\[
W=\operatorname{span}\{v_1,\ldots,v_{k+1}\}.
\]

Because these vectors are orthonormal,

\[
\dim(W)=k+1.
\]

\textbf{Step 2: Find a unit vector
that $B$ sends to zero.}

Restrict the transformation $B$ to $W$:

\[
B|_W:W\longrightarrow\mathbb{R}^m.
\]

Its rank cannot exceed $k$.
By the rank-nullity theorem,

\[
\begin{aligned}
\dim\ker(B|_W)
&=\dim(W)-\operatorname{rank}(B|_W)\\
&\geq(k+1)-k=1.
\end{aligned}
\]

Therefore, there exists a unit vector
$x\in W$ such that

\[
\boxed{Bx=0,\qquad \|x\|_2=1.}
\]

This is the crucial observation:
any matrix of rank at most $k$
necessarily loses at least one
nonzero direction from the
$(k+1)$-dimensional space $W$.

\textbf{Step 3: Estimate the error
on this special vector.}

Since $x\in W$, write

\[
x=\sum_{i=1}^{k+1}c_i v_i,
\qquad
\sum_{i=1}^{k+1}c_i^2=1.
\]

Using $Av_i=\sigma_i u_i$, we obtain

\[
Ax=\sum_{i=1}^{k+1}c_i\sigma_i u_i.
\]

The left singular vectors are
orthonormal, so

\[
\begin{aligned}
\|Ax\|_2^2
&=\sum_{i=1}^{k+1}c_i^2\sigma_i^2\\
&\geq \sigma_{k+1}^2
\sum_{i=1}^{k+1}c_i^2\\
&=\sigma_{k+1}^2.
\end{aligned}
\]

But $Bx=0$, and hence

\[
\|(A-B)x\|_2
=\|Ax\|_2\geq\sigma_{k+1}.
\]

By the definition of the spectral norm,

\[
\boxed{
\|A-B\|_2
=\max_{\|y\|_2=1}\|(A-B)y\|_2
\geq\sigma_{k+1}.
}
\]

Thus, no rank-$k$ matrix can
achieve a smaller spectral-norm error.

\textbf{Step 4: Show that $A_k$
achieves the lower bound.}

Recall that

\[
A-A_k=\sum_{i=k+1}^{r}\sigma_i u_i v_i^T.
\]

This residual has singular values
$\sigma_{k+1},\ldots,\sigma_r$.
Its largest singular value is
$\sigma_{k+1}$. Therefore,

\[
\boxed{\|A-A_k\|_2=\sigma_{k+1}.}
\]

Combining the lower bound with
this equality yields

\[
\boxed{
\min_{\operatorname{rank}(B)\leq k}
\|A-B\|_2=\sigma_{k+1}.
}
\]

\end{proof}

% -------------------------------------------------
\subsection{What Do These Results Tell Us About PCA?}
% -------------------------------------------------

Let us now return to the centered data
matrix used in Principal Component Analysis.

Its truncated SVD satisfies

\[
\begin{aligned}
A_k
&=\sum_{i=1}^{k}\sigma_i u_i v_i^T\\
&=U_kU_k^TA.
\end{aligned}
\]

Thus, the matrix obtained by
projecting onto the principal subspace
and reconstructing the observations
is exactly the truncated SVD.

The Eckart--Young theorem now supplies
the promised global optimality statement:

\[
\boxed{
\begin{aligned}
\|A-A_k\|_F
&\leq\|A-B\|_F,\\
\|A-A_k\|_2
&\leq\|A-B\|_2
\end{aligned}
}
\]

for every $B\in\mathbb{R}^{m\times n}$
with $\operatorname{rank}(B)\leq k$.

The Frobenius result minimizes total
squared reconstruction error, while
the spectral result minimizes the
worst directional reconstruction error.

The theorem does not require the competing
approximation $B$ to be an orthogonal
projection.

There is also a subtle point about
uniqueness. The theorem guarantees
that the truncated SVD is optimal,
but does not guarantee that it is
the only optimal approximation.

In particular, repeated singular values
may lead to different optimal choices, which we may later discuss in another note.

We have therefore connected the
statistical objective of preserving
variance to the linear-algebraic
objective of optimal low-rank
approximation.

What began as a geometric intuition
is now supported by a rigorous theorem.

% =================================================
\section{Numerical Implementation of PCA Using Python}
% =================================================

Having established the mathematical foundations
of Principal Component Analysis and proved the
Eckart--Young theorem, we now illustrate these
results through numerical experiments using Python.

We implement PCA directly using NumPy's Singular
Value Decomposition rather than relying on a
prebuilt PCA algorithm.

Throughout our implementation, we follow the
convention adopted in this note:

\[
X\in\mathbb{R}^{m\times n},
\]

where rows represent features and columns
represent observations.

\subsection{Implementation Using a Small Dataset}

We begin with a simple dataset containing
two features and four observations:

\[
X=
\begin{bmatrix}
8 & 9 & 11 & 12\\
4 & 4 & 6 & 6
\end{bmatrix}.
\]

We first center the data by subtracting
the mean of each feature. We then compute
the covariance matrix and the Singular
Value Decomposition of the centered data.

Retaining only the first principal direction,
we reduce the original two-dimensional
observations to one-dimensional coordinates.

The numerical experiment demonstrates that
the first principal component preserves
approximately $97.92\%$ of the total variance.

We also reconstruct the observations and
verify that the squared Frobenius error and
spectral error agree with the formulas
established by the Eckart--Young theorem.

\subsection{Implementation Using a Real Dataset}

Next, we consider a larger dataset consisting
of handwritten digit images.

The dataset contains $1797$ observations,
each represented by $64$ numerical features
corresponding to the pixels of an $8\times8$
image.

Following our convention, the data matrix
has dimensions

\[
X\in\mathbb{R}^{64\times1797}.
\]

After centering the observations, we compute
the SVD and select the smallest number of
principal components required to preserve
at least $90\%$ of the total variance.

In this experiment, $21$ principal components
preserve approximately $90.32\%$ of the
total variance.

Thus, we reduce the original $64$ features
to only $21$ principal component coordinates.

We subsequently reconstruct the original
images and compare them with their
low-rank approximations.

Finally, we verify numerically that

\[
\|A-A_k\|_F^2
=
\sum_{i=k+1}^{r}\sigma_i^2
\]

and

\[
\|A-A_k\|_2=\sigma_{k+1},
\]

confirming the error identities established
by the Eckart--Young theorem.

\subsection{Python Source Code}

The complete Python implementation for both
experiments is provided below.

The program also produces graphical illustrations
of the principal directions, explained variance,
and reconstructed observations.

\lstinputlisting[
    language=Python,
    caption={Numerical implementation of PCA and
    verification of the Eckart--Young theorem},
    label={lst:pca_python},
    basicstyle=\scriptsize\ttfamily,
    breaklines=true,
    columns=fullflexible,
    keepspaces=true,
    showstringspaces=false,
    numbers=left,
    numberstyle=\tiny
]{PCA_implementations.py}


% =================================================
\section{Conclusion and Future Directions}
% =================================================

In this note, we have attempted to develop
Principal Component Analysis from its
fundamental mathematical principles,
rather than treating it merely as a
computational procedure.

We began with a brief review of Singular
Value Decomposition, examining its
relationship with the row space, column
space, singular vectors, and orthogonal
projections. We then introduced the
geometric motivation behind PCA and
explained how variance and covariance
allow us to identify the dominant
directions of variation in a dataset.

Building upon these ideas, we developed
the PCA procedure step by step:
centering the observations, computing
the singular value decomposition,
selecting the principal directions,
projecting the original observations,
and reconstructing their low-dimensional
approximations.

An important objective of our discussion
was to explain not only how these
operations are performed, but also
why they are mathematically justified.

In particular, we established that
the first $k$ left singular vectors
identify a subspace that maximizes
the total retained variance among all
$k$-dimensional orthogonal projections.

We subsequently showed that maximizing
the retained variance is equivalent
to minimizing the squared Frobenius
reconstruction error.

Finally, through the Eckart--Young theorem,
we established a considerably stronger
result: the truncated Singular Value
Decomposition provides an optimal
rank-$k$ approximation among all matrices
of rank at most $k$, under both the
Frobenius and spectral norms.

Thus, the approximation

\[
\boxed{
A_k=\sum_{i=1}^{k}
\sigma_i u_i v_i^T
}
\]

is not merely a convenient construction.
Its optimality follows from a rigorous
mathematical theorem.

We also implemented PCA using Python,
first on a small, manually constructed
dataset and subsequently on a larger
dataset of handwritten digit images.
These experiments demonstrated how
dimensionality reduction and reconstruction
work in practice, while numerically
verifying the error formulas predicted
by the Eckart--Young theorem.

Nevertheless, despite its mathematical
elegance and widespread usefulness,
PCA is far from a perfect method,
and the practical use of SVD is not
without limitations.

One important difficulty concerns
categorical and mixed-type datasets.
Ordinary PCA is based on numerical
coordinates, Euclidean distances,
and variance. When categorical
variables are represented by
arbitrary numerical labels, these
operations may introduce artificial
relationships between categories.

Even when categorical observations
are converted into numerical vectors
through methods such as one-hot
encoding, the resulting Euclidean
geometry may not adequately reflect
the underlying structure of the data.

It is important to emphasize that
this is primarily a limitation of
ordinary PCA and the way the data
is represented, rather than a failure
of the SVD theorem itself.

Methods such as Multiple Correspondence
Analysis and techniques designed
specifically for mixed numerical and
categorical data provide alternative
approaches in such situations.

Another important challenge arises
in large-scale scientific computing.

Although highly reliable numerical
algorithms for computing SVD already
exist, calculating a full decomposition
of an extremely large matrix can
require substantial computational
time and memory.

This becomes particularly significant
in applications involving high-resolution
images, large scientific simulations,
massive datasets, or continuously
arriving observations.

Moreover, PCA identifies linear
subspaces and maximizes variance,
which may not capture nonlinear
relationships or necessarily preserve
the information most relevant to a
particular predictive task.

Fortunately, these limitations have
also motivated considerable progress
in numerical linear algebra and
machine learning.

Research on randomized and truncated
SVD algorithms, matrix sketching,
incremental PCA, distributed computing,
and hardware-accelerated numerical
methods continues to improve the
scalability of low-rank approximation.

At the same time, nonlinear
dimensionality-reduction techniques,
methods for categorical and mixed
data, and more flexible representation
models offer promising alternatives
when the assumptions of classical
PCA are inappropriate.

These developments suggest that the
future of dimensionality reduction
will not necessarily consist of
replacing SVD, but rather of extending
its computational reach and combining
its mathematical strengths with methods
better suited to increasingly complex
forms of data.

We may therefore expect further
improvements in the efficiency,
scalability, and applicability
of low-dimensional approximation
methods in the coming years.

We conclude by emphasizing that
understanding the limitations of
a mathematical method is just as
important as understanding its
strengths.

The Eckart--Young theorem tells us
precisely when and why truncated
SVD is optimal. It does not claim
that low-rank approximation is
the appropriate objective for
every possible dataset.

Recognizing this distinction allows
us to appreciate both the remarkable
mathematical beauty of PCA and the
important challenges that remain.

Our hope is that this exposition
has helped the reader see PCA not
as an obscure black-box algorithm,
but as a natural and elegant
consequence of elementary linear
algebra, geometry, and optimization,
while also encouraging further
exploration of the rapidly evolving
field of numerical data analysis.


% =================================================
% REFERENCES
% =================================================

\begin{thebibliography}{99}

\bibitem{strang2006}
G. Strang,
\textit{Linear Algebra and Its Applications},
4th ed.,
Cengage Learning, 2006.

\bibitem{kutz2013}
J. N. Kutz,
\textit{Data-Driven Modeling \& Scientific
Computation: Methods for Complex Systems
\& Big Data},
Oxford University Press, 2013.

\bibitem{belhaizaPCA}
S. J. Belhaiza,
\textit{Principal Component Analysis},
MATH 557 video lecture,
King Fahd University of Petroleum
and Minerals (KFUPM).
Available online:
\url{https://www.youtube.com/watch?v=lpBCHV1rQTI}.

\bibitem{jolliffe2002}
I. T. Jolliffe,
\textit{Principal Component Analysis},
2nd ed.,
Springer Series in Statistics,
Springer, New York, 2002.
\url{https://doi.org/10.1007/b98835}.

\bibitem{trefethen1997}
L. N. Trefethen and D. Bau III,
\textit{Numerical Linear Algebra},
Society for Industrial and
Applied Mathematics (SIAM), 1997.

\bibitem{golub2013}
G. H. Golub and C. F. Van Loan,
\textit{Matrix Computations},
4th ed.,
Johns Hopkins University Press, 2013.

\bibitem{eckart1936}
C. Eckart and G. Young,
``The Approximation of One Matrix
by Another of Lower Rank,''
\textit{Psychometrika},
vol. 1, no. 3,
pp. 211--218, 1936.
\url{https://doi.org/10.1007/BF02288367}.

\bibitem{mirsky1960}
L. Mirsky,
``Symmetric Gauge Functions
and Unitarily Invariant Norms,''
\textit{The Quarterly Journal
of Mathematics},
vol. 11, no. 1,
pp. 50--59, 1960.
\url{https://doi.org/10.1093/qmath/11.1.50}.

\bibitem{halko2011}
N. Halko, P. G. Martinsson,
and J. A. Tropp,
``Finding Structure with Randomness:
Probabilistic Algorithms for Constructing
Approximate Matrix Decompositions,''
\textit{SIAM Review},
vol. 53, no. 2,
pp. 217--288, 2011.
\url{https://doi.org/10.1137/090771806}.

\end{thebibliography}


\end{document}